\documentclass[11pt]{article}
\usepackage[a4paper,margin=22mm]{geometry}
\usepackage{fontspec}
\setmainfont{LinLibertine_R.otf}[BoldFont=LinLibertine_RB.otf,ItalicFont=LinLibertine_RI.otf,BoldItalicFont=LinLibertine_RBI.otf]
\usepackage{xcolor}
\usepackage{hyperref}
\hypersetup{colorlinks=true,urlcolor=blue}
\setlength{\parindent}{0pt}
\setlength{\parskip}{6pt}
\sloppy
\emergencystretch=2em
\begin{document}
\section*{Historical Britannica: extracted source}
Encyclopaedia Britannica - Sixth edition, Volume 13, MAT-MIC - EB.11\par
{\small Source row 2074. Parser: nvidia/nemotron-3.5-parse.\par}
{\small This is model-extracted text, not a corrected transcription. OCR errors, reading order and region alignment require human review.\par}
\hrule\medskip
{\small\color{red} Rendering note: glyphs unsupported by this document font appear as boxed Unicode codes. Original characters remain in the raw model output and HTML transcription. No corrected reading is inferred.\par}

{\small\itshape MECHANICS.\par}

{\scriptsize\color{gray} Region region-0; provider order 0.\par}
\medskip

{\small\itshape 103\par}

{\scriptsize\color{gray} Region region-1; provider order 1.\par}
\medskip

heer, d gives \textbackslash{}(\textbackslash{}frac\{d\textasciicircum{}2\textbackslash{}pi xd\}\{2\}\textbackslash{}) for the sum of all its particles. But as this is to be multiplied by the sum of the squares of all the distances of the particles from the centre C, we shall have \textbackslash{}(fp r\textasciicircum{}r=a\textasciicircum{}\{t+x\}x\textbackslash{}lambda d\textbackslash{}) . But the number of particles in the cylinder, or the mass \textbackslash{}(\textbackslash{}mu\textbackslash{}) of the cy- linder, is \textbackslash{}(a\textasciicircum{}\textbackslash{}alpha\textbackslash{}pi\textbackslash{}lambda d\textbackslash{}) , therefore, substituting \textbackslash{}(\textbackslash{}mu\textbackslash{}), instead of this value of it in the preceding equation, we have \textbackslash{}(fpr\textasciicircum{}r=\textbackslash{}frac\{\textbackslash{}mu a\textasciicircum{}2\}\{2\}\textbackslash{}), and, dividing both sides by \_n\_, we have
\par

{\scriptsize\color{gray} Region region-2; provider order 2.\par}
\medskip

\textbackslash{}(\textbackslash{}int\_\{pr\textasciicircum{}r\}\textasciicircum{}n\}\{\textbackslash{}nu\}=\textbackslash{}frac\{\textbackslash{}mu a\textasciicircum{}2\}\{2n\}\textbackslash{}) and, extracting the square root and mul- tiplying by 180, it becomes
\par

{\scriptsize\color{gray} Region region-3; provider order 3.\par}
\medskip

\textbackslash{}(\textbackslash{}sqrt[\textbackslash{}rho r\textasciicircum{}r]\textbackslash{}big|\textbackslash{}xi\textbackslash{}times 18o\textasciicircum{}s=\textbackslash{}frac\{a\textasciicircum{}1\}\{2n\}\textbackslash{}big|\textbackslash{}xi\textbackslash{}times 18o\textasciicircum{}s\textbackslash{}). Therefore
\par

{\scriptsize\color{gray} Region region-4; provider order 4.\par}
\medskip

\textbackslash{}(T=\textbackslash{}frac\{\textbackslash{}mu a\textasciicircum{}2\}\{2n\}\textbackslash{}big|\textbackslash{}xi\textbackslash{}times 18o\textbackslash{}), and since \textbackslash{}(\textbackslash{}int\_rp\textasciicircum{}\{pr\textasciicircum{}r\}=l/g\textbackslash{}), \textbackslash{}(\textbackslash{}frac\{i\}\{g\}\textbackslash{}), \textbackslash{}(\textbackslash{}frac\{\textbackslash{}mu a\textasciicircum{}2\}\{2n\}=\textbackslash{}) \_l g\_, and by reduction \textbackslash{}(\textbackslash{}frac\{v-\textbackslash{}sum\textbackslash{}mu q\}\{2L\}\textbackslash{}). But \_g\_ as is the weight W of the cylinder, therefore, by substituting W instead of \_g\_ \textbackslash{}(\textbackslash{}mu\textbackslash{}), we obtain \_n\_=P<sup>\_a\_</sup>/2L, very simple formu
\par

{\scriptsize\color{gray} Region region-5; provider order 5.\par}
\medskip

Ia for determining the value of n from experiments.
\par

{\scriptsize\color{gray} Region region-6; provider order 6.\par}
\medskip

If it were required to find a weight Q, which, acting at the extremity of a lever I, would have no momentum equal to the momentum of the force of torsion when the angle of torsion is A--M, we must make Q\fbox{\texttt{U+2A09}}L=nX A--M.
\par

{\scriptsize\color{gray} Region region-7; provider order 7.\par}
\medskip

346. In the preceding investigation we have suppo- soil, what is conformable to systemically, that the forces of torsion is proportional to the angle of torsion, which gives us \_n\_xxA--M for the momentum of that force. Let us now suppose that this momentum is altered by any quantity S, then the momentum of the force of torsion will become \_n\_xxA--M--S’, and the general equation will assume this form
\par

{\scriptsize\color{gray} Region region-8; provider order 8.\par}
\medskip

nXA--M--S\fbox{\texttt{U+2573}}:=u \textbackslash{}(\textbackslash{}int(p r\textasciicircum{}a\textasciicircum{}a)\textbackslash{}),
\par

{\scriptsize\color{gray} Region region-9; provider order 9.\par}
\medskip

and by multiplying in place of \_t\_ its value \textbackslash{}(\textbackslash{}frac\{oM\}\{a\}\textbackslash{}), and
\par

{\scriptsize\color{gray} Region region-10; provider order 10.\par}
\medskip

taking the flient, we have
\par

{\scriptsize\color{gray} Region region-11; provider order 11.\par}
\medskip

\textbackslash{}(n\textbackslash{}times 2\textbackslash{}Lambda Mm-M\textasciicircum{}\{-2\}\textbackslash{}int 8M\textbackslash{}equiv d'\textbackslash{}frac\{d'r\textasciicircum{}a\}\{a\textasciicircum{}k\}\textbackslash{})
\par

{\scriptsize\color{gray} Region region-12; provider order 12.\par}
\medskip

Now, in other to find the value of P or a complete oscillation, we must divide the oscillation into two parts, the first from B to A, where the force of torsion accelerates the velocity \_u\_, while the retarded force, arising from the resistance of the air and the imper- flection of elasticity, diminishes the velocity \_u\_, and the
\par

{\scriptsize\color{gray} Region region-13; provider order 13.\par}
\medskip

second from A to B’, where the force of torsion, as well as the other forces, concerner in diminishing \_n\_ or retarding the motion.
\par

{\scriptsize\color{gray} Region region-14; provider order 14.\par}
\medskip

347. \_Ex.\_ 1. If S=m×A--M’, we shall have for the state of motion in the first portion BA
\par

{\scriptsize\color{gray} Region region-15; provider order 15.\par}
\medskip

\textbackslash{}(n\textbackslash{}times 2\textbackslash{}Lambda M-M+\textbackslash{}frac\{2m\textbackslash{}lambda A-M\textasciicircum{}\{+1\}\}\{+1+\}\textbackslash{}frac\{zm\textbackslash{}lambda\textasciicircum{}+1\}\{+1+\}=\textbackslash{}omega t\textbackslash{}int p\textasciicircum{}a\textasciicircum{}a\textasciicircum{}a\textbackslash{})
\par

{\scriptsize\color{gray} Region region-16; provider order 16.\par}
\medskip

Hence, when the angle of torsion becomes equal to no- thing, or A--M=0, we have
\par

{\scriptsize\color{gray} Region region-17; provider order 17.\par}
\medskip

\textbackslash{}(n\textbackslash{}Delta\textasciicircum{}+\textbackslash{}frac\{2m\textbackslash{}lambda\textasciicircum{}+1\}\{++1\}=\textbackslash{}UU\textbackslash{}int p\textasciicircum{}a\textasciicircum{}a\textasciicircum{}a\textasciicircum{}a,\textbackslash{})
\par

{\scriptsize\color{gray} Region region-18; provider order 18.\par}
\medskip

which dividing by \textbackslash{}(\textbackslash{}int\textasciicircum{}p\textasciicircum{}p\textasciicircum{}a\textbackslash{}), becomes
\par

{\scriptsize\color{gray} Region region-19; provider order 19.\par}
\medskip

\textbackslash{}(U!=\textbackslash{}frac\{n\textbackslash{}Delta\textasciicircum{}+\textbackslash{}frac\{2m\textbackslash{}lambda\textasciicircum{}+1\}\{++1\}\}\textbackslash{})
\par

{\scriptsize\color{gray} Region region-20; provider order 20.\par}
\medskip

Let us now consider the other part of the motion from A to B’, and suppose the angle AC U=M, we shall find, by calling U the velocity of the point A,
\par

{\scriptsize\color{gray} Region region-21; provider order 21.\par}
\medskip

\textbackslash{}(nMP\textasciicircum{}a\}\{2\}++\textbackslash{}frac\{M\textasciicircum{}kr\textasciicircum{}+r+1\}=\textbackslash{}U\textasciicircum{}\{-a\textasciicircum{}*\}/\textbackslash{}sum\textbackslash{}int\textasciicircum{}p\textasciicircum{}p\textasciicircum{}a\textasciicircum{}r.\textbackslash{})
\par

{\scriptsize\color{gray} Region region-22; provider order 22.\par}
\medskip

Then, by substituting instead of U its value as lately found, and taking the floants, we shall have, when the velocity vanishes, or when the oscillation is finished,
\par

{\scriptsize\color{gray} Region region-23; provider order 23.\par}
\medskip

\textbackslash{}(\textbackslash{}Delta-M=\textbackslash{}frac\{2m\}\{nx++2\}\textbackslash{}times\textbackslash{}frac\{A\textasciicircum{}\{+++\}M\textasciicircum{}\{+1\}\}\{+1+\},\textbackslash{})
\par

{\scriptsize\color{gray} Region region-24; provider order 24.\par}
\medskip

and if the retarding forces are such, that at each oscilla- tion, the amplitude is a little diminished, we shall have for the approximate value of A--M
\par

{\scriptsize\color{gray} Region region-25; provider order 25.\par}
\medskip

\textbackslash{}(\textbackslash{}Lambda-M=\textbackslash{}frac\{2m\textbackslash{}lambda\}\{n\textbackslash{}px++i\}\textbackslash{}),
\par

{\scriptsize\color{gray} Region region-26; provider order 26.\par}
\medskip

and if the angle A--M is so small that it may be treated as a common flatened quantity, we shall then have for any number of oscillations
\par

{\scriptsize\color{gray} Region region-27; provider order 27.\par}
\medskip

\textbackslash{}(N\textbackslash{}times\textbackslash{}frac\{2m\}\{nx++1\}=\textbackslash{}frac\{1\}\{1+\}\textbackslash{}times\textbackslash{}frac\{1\}\{M\textasciicircum{}\{+-i\}\}+\textbackslash{}frac\{1\}\{A\textasciicircum{}\{+1\}\}.\textbackslash{})
\par

{\scriptsize\color{gray} Region region-28; provider order 28.\par}
\medskip

where M represents the angle to which A becomes equal after any number of oscillations N. Hence we obtain
\par

{\scriptsize\color{gray} Region region-29; provider order 29.\par}
\medskip

\textbackslash{}(M=\textbackslash{}frac\{1\}\{(N\textbackslash{}times\textbackslash{}frac\{2m\textbackslash{}lambda+1\}\{n\textbackslash{}times+1\}+\textbackslash{}frac\{1\}\{A\textasciicircum{}\{+1\}\})\textbackslash{}times\textbackslash{}frac\{1\}\{r+1\}\textbackslash{}),
\par

{\scriptsize\color{gray} Region region-30; provider order 30.\par}
\medskip

which determines the value of M after any number of oscillations N.
\par

{\scriptsize\color{gray} Region region-31; provider order 31.\par}
\medskip

348. \_Ex.\_ 2. If S=m×A--M’,+m×A--Mp, m’ and \_r\_ being different values of m and \_r\_, we shall obtain by following the mode of investigation in the last ex- ample,
\par

{\scriptsize\color{gray} Region region-32; provider order 32.\par}
\medskip

\textbackslash{}(n\textbackslash{}times\textbackslash{}Delta-M=\textbackslash{}frac\{2m\}\{r+1\}\textbackslash{}times\textbackslash{}frac\{A\textasciicircum{}\{++\}+M\textasciicircum{}\{+1\}\}\{++m\}/\textbackslash{}times\textbackslash{}frac\{A\textasciicircum{}\{++\}xM\textasciicircum{}\{+1\}\}\{A+M\}.\textbackslash{})
\par

{\scriptsize\color{gray} Region region-33; provider order 33.\par}
\medskip

and if the retarding force is much less than the force of torsion, we shall have for an approximate value of
\par

{\scriptsize\color{gray} Region region-34; provider order 34.\par}
\medskip

\textbackslash{}(n\textbackslash{}times\textbackslash{}Delta-M.\textbackslash{})
\par

{\scriptsize\color{gray} Region region-35; provider order 35.\par}
\medskip
\newpage\section*{Source record}
bri-2074-d43e4278eb9d\par
Image SHA-256: \texttt{d43e4278eb9dfae4b0b0c99f2735aa0e}\par\texttt{223dbe518a7ee9fb465cf46135364a04}\par
Dataset revision: \texttt{309f765b1ecc83e0a7f7}\texttt{762e44b4a927a7555443}\par
National Library of Scotland, Encyclopaedia Britannica Illustrated (CC0).\par
\url{https://huggingface.co/datasets/NationalLibraryOfScotland/encyclopaedia_britannica_illustrated/viewer/default/train?row=2074}
\end{document}
